% GENERATED FILE. Edit canonical lessons, then run npm run generate:books.
% canonical-source-hash: fnv1a64:cfbb8688313ec153
% canonical-lessons: ZH-C43-qu, ZH-W43-qu, ZH-C43-hui, ZH-C43-zou, ZH-C43-chumen, ZH-C43-xuexiao

\chapter{Coming and Going}
\label{ch:zh-coming-and-going}
\hlchaptermodality{43}

\begin{chapteropening}
\textbf{By the end of this chapter:} \emph{I can say where I am going and that I am going home.}

Complete the last lesson of chapter 43: I can say where I am going and that I am going home.
\end{chapteropening}

\section[\texorpdfstring{\zh{去} (qù)}{去 (qù)}]{\zh{去} (qù) — to go}

\label{lesson:ZH-C43-qu}



\begin{quote}
Before the new one: say the Mandarin for cheap, then the Mandarin for expensive.
\end{quote}

\subsection*{You'll want to know: \zh{去}}
\textbf{\zh{去}} — \emph{qù} — \textquotedblleft{}to go\textquotedblright{}.

\begin{quote}
\textbf{\zh{我去大学}} \emph{wǒ qù dàxué} — I'm going to the university.
\end{quote}

The new character, in pieces: \zh{去} — a shape like \zh{十} standing on a long bar, and under it a two-stroke piece. The next lesson puts it in your hand.

\begin{cousinweb}
\textbf{\zh{去}} (\emph{qù}) is inherited from Middle Chinese \emph{khjoH} (Guangyun, in Baxter's transcription), glossed 'to leave'.
\end{cousinweb}

\begin{grammarlens}[title={What you've built}]
One more word for everyday answers.
\end{grammarlens}

\subsection*{Your turn}
\begin{itemize}
\raggedright
  \item \emph{Say it:} \emph{qù}
  \item \emph{Say it:} \emph{qù}, once more
  \item \emph{Say it:} \emph{wǒ qù dàxué} — I'm going to the university.
  \item \emph{Recall:} say the Mandarin for cheap, then the Mandarin for expensive, then say \emph{qù} again
\end{itemize}

\begin{tcolorbox}[breakable,colback=teal!4,colframe=teal!35!black,title={Before you move on}]
What does \zh{去} mean? (\textquotedblleft{}To go\textquotedblright{}.) Say it once more.
\end{tcolorbox}
\section[\texorpdfstring{\zh{去} (qù)}{去 (qù)}]{\zh{去} — a cross on a bar, and a corner below}

\label{lesson:ZH-W43-qu}



\begin{quote}
The word first, then the character inside it.

\begin{itemize}
\raggedright
  \item \emph{Recall:} say \emph{qù} — \textbf{R1}, one lesson back, before it decays
\end{itemize}

\begin{quote}
\zh{十} \zh{大}
\end{quote}

\emph{Read it:} each one aloud, with nothing said first — \emph{shí}, \emph{dà}
\end{quote}

\subsection*{Script — a cross on a bar, and a corner below}
\hlblockfigure{\detokenize{figures/ZH-W43-qu-filmstrip.pdf}}{How \zh{去} is written, stroke by stroke}

\begin{quote}
\zh{去}
\end{quote}

\textbf{Two pieces, top before bottom.}

\textbf{The top piece, three strokes.} A short horizontal, a vertical through it, then a long horizontal under them: \zh{十} standing on a wide base.

\textbf{The bottom piece, two strokes.} A stroke that falls to the left and turns out to the right, then a dot inside the turn.

3 + 2 = five.

\subsection*{Your turn}
\begin{itemize}
\raggedright
  \item \emph{Write it:} \zh{去}, all five, with the model in front of you
  \item \emph{Cover:} write it again from memory, and count the pieces, not the strokes
\end{itemize}

\begin{tcolorbox}[breakable,colback=teal!4,colframe=teal!35!black,title={Before you move on}]
How many pieces in \zh{去}? (\textbf{Two}.) Which character you know is in the top piece? (\textbf{\zh{十}}, on a long bar.) What is the last stroke? (\textbf{A dot}.)

Source: \href{https://raw.githubusercontent.com/chanind/hanzi-writer-data/68d10a4b21150cae5e1ebbd223eed289cf32d90c/data/\%E5\%8E\%BB.json}{Hanzi Writer Data \zh{去}, pinned stroke paths and medians}.
\end{tcolorbox}
\section[huí]{huí — to go back; to return}

\label{lesson:ZH-C43-hui}



\begin{quote}
Before the new one: say the Mandarin for expensive, then the Mandarin for to go.
\end{quote}

\subsection*{You'll want to know: huí}
\emph{huí} — \textquotedblleft{}to go back; to return\textquotedblright{}.

\begin{quote}
\emph{wǒ huí jiā} — I'm going home.
\end{quote}

Its characters come later; for now it is a word for the ear.

\begin{cousinweb}
\emph{huí} is inherited from Middle Chinese \emph{hwoj} (Guangyun, in Baxter's transcription), glossed 'to turn'.
\end{cousinweb}

\begin{grammarlens}[title={What you've built}]
One more word for everyday answers.
\end{grammarlens}

\subsection*{Your turn}
\begin{itemize}
\raggedright
  \item \emph{Say it:} \emph{huí}
  \item \emph{Say it:} \emph{huí}, once more
  \item \emph{Say it:} \emph{wǒ huí jiā} — I'm going home.
  \item \emph{Recall:} say the Mandarin for expensive, then the Mandarin for to go, then say \emph{huí} again
\end{itemize}

\begin{tcolorbox}[breakable,colback=teal!4,colframe=teal!35!black,title={Before you move on}]
What does huí mean? (\textquotedblleft{}To go back; to return\textquotedblright{}.) Say it once more.
\end{tcolorbox}
\section[zǒu]{zǒu — to walk; to leave}

\label{lesson:ZH-C43-zou}



\begin{quote}
Before the new one: say the Mandarin for to go, then the Mandarin for to go back.
\end{quote}

\subsection*{You'll want to know: zǒu}
\emph{zǒu} — \textquotedblleft{}to walk; to leave\textquotedblright{}.

\begin{quote}
\emph{wǒ zǒu le} — I'm off.
\end{quote}

Its characters come later; for now it is a word for the ear.

\begin{cousinweb}
\emph{zǒu} is inherited from Middle Chinese \emph{tsuwX} (Guangyun, in Baxter's transcription), glossed 'to hurry'. The sense 'to walk' is not given there, so when it arrived is not settled.
\end{cousinweb}

\begin{grammarlens}[title={What you've built}]
One more word for everyday answers.
\end{grammarlens}

\subsection*{Your turn}
\begin{itemize}
\raggedright
  \item \emph{Say it:} \emph{zǒu}
  \item \emph{Say it:} \emph{zǒu}, once more
  \item \emph{Say it:} \emph{wǒ zǒu le} — I'm off.
  \item \emph{Recall:} say the Mandarin for to go, then the Mandarin for to go back, then say \emph{zǒu} again
\end{itemize}

\begin{tcolorbox}[breakable,colback=teal!4,colframe=teal!35!black,title={Before you move on}]
What does zǒu mean? (\textquotedblleft{}To walk; to leave\textquotedblright{}.) Say it once more.
\end{tcolorbox}
\section[chūmén]{chūmén — to go out}

\label{lesson:ZH-C43-chumen}



\begin{quote}
Before the new one: say the Mandarin for to go back, then the Mandarin for to walk.
\end{quote}

\subsection*{You'll want to know: chūmén}
\emph{chūmén} — \textquotedblleft{}to go out\textquotedblright{}.

\begin{quote}
\emph{māma chūmén le} — Mum's gone out.
\end{quote}

Its characters come later; for now it is a word for the ear.

\begin{cousinweb}
\emph{chūmén} is built from chū 'to come forth' and mén 'door', both inherited from Middle Chinese, \emph{tsyhwit} and \emph{mwon} (Guangyun, in Baxter's transcription).
\end{cousinweb}

\begin{grammarlens}[title={What you've built}]
One more word for everyday answers.
\end{grammarlens}

\subsection*{Your turn}
\begin{itemize}
\raggedright
  \item \emph{Say it:} \emph{chūmén}
  \item \emph{Say it:} \emph{chūmén}, once more
  \item \emph{Say it:} \emph{māma chūmén le} — Mum's gone out.
  \item \emph{Recall:} say the Mandarin for to go back, then the Mandarin for to walk, then say \emph{chūmén} again
\end{itemize}

\begin{tcolorbox}[breakable,colback=teal!4,colframe=teal!35!black,title={Before you move on}]
What does chūmén mean? (\textquotedblleft{}To go out\textquotedblright{}.) Say it once more.
\end{tcolorbox}
\section[xuéxiào]{xuéxiào — school}

\label{lesson:ZH-C43-xuexiao}



\begin{quote}
Before the new one: say the Mandarin for to walk, then the Mandarin for to go out.
\end{quote}

\subsection*{You'll want to know: xuéxiào}
\emph{xuéxiào} — \textquotedblleft{}school\textquotedblright{}.

\begin{quote}
\emph{wǒmen qù xuéxiào} — We're off to school.
\end{quote}

Its characters come later; for now it is a word for the ear.

\begin{cousinweb}
\emph{xuéxiào} is built from \zh{学} 'learning', inherited from Middle Chinese \emph{haewk} (Guangyun, in Baxter's transcription), and a second character the dictionary records as \emph{haewH}, the title of an officer. The school sense is not given there, so how it came to mean school is not settled.
\end{cousinweb}

\begin{grammarlens}[title={What you've built}]
One more word for everyday answers.
\end{grammarlens}

\subsection*{Your turn}
\begin{itemize}
\raggedright
  \item \emph{Say it:} \emph{xuéxiào}
  \item \emph{Say it:} \emph{xuéxiào}, once more
  \item \emph{Say it:} \emph{wǒmen qù xuéxiào} — We're off to school.
  \item \emph{Recall:} say the Mandarin for to walk, then the Mandarin for to go out, then say \emph{xuéxiào} again
\end{itemize}

\begin{tcolorbox}[breakable,colback=teal!4,colframe=teal!35!black,title={Before you move on}]
What does xuéxiào mean? (\textquotedblleft{}School\textquotedblright{}.) Say it once more.
\end{tcolorbox}
